\chapter{Iron Deficiency}
\index{Iron Deficiency}
\label{sec:Iron Deficiency}

\section{Overview}
Iron deficiency is the most common nutritional deficiency worldwide and the leading cause of anemia, affecting approximately 25\% of the global population. This metabolic disorder involves dysregulation of normal bodily processes, often requiring ongoing monitoring and management. It is increasingly common due to lifestyle and dietary factors. Metabolic disorders involve disruption of normal biochemical processes essential for health. These conditions may result from enzyme deficiencies, hormonal imbalances, or lifestyle factors. Many metabolic disorders are chronic and require ongoing management. Lifestyle modifications are often the cornerstone of treatment.
\section{Causes}
The underlying mechanisms involve complex interactions between genetic predisposition and environmental factors. Disruption of normal physiological processes leads to the characteristic features of the condition. Multiple contributing factors may act synergistically to trigger or worsen the condition.

\begin{itemize}[noitemsep]
  \item This condition happens because of inadequate dietary iron intake, poor bioavailability (non-heme iron from plant sources), increased demands (pregnancy, growth, menstruation), and chronic blood loss (digestive tract bleeding, menorrhagia, hookworm)
  \item Iron is essential for hemoglobin synthesis, myoglobin, and numerous enzymes involved in electron transport and DNA synthesis.
\end{itemize}
\section{Risk Factors}
\begin{itemize}[noitemsep]
  \item Genetic predisposition and family history
  \item Advancing age
  \item Lifestyle factors (diet, exercise, smoking, alcohol)
  \item Environmental exposures
  \item Underlying medical conditions
  \item Medication use
\end{itemize}
\section{Signs and Symptoms}

\subsection{Common symptoms}
\begin{itemize}[noitemsep]
  \item Clinical features manifest as microcytic hypochromic anemia including fatigue, pallor, weakness, shortness of breath on exertion, tachycardia, headache, and epithelial changes including koilonychia (spoon nails), glossitis, angular stomatitis, and pagophagia (pica for ice).
  \item Pain or discomfort in the affected area
  \end{itemize}

\subsection{Emergency warning signs}
\begin{itemize}[noitemsep]
  \item Severe or worsening symptoms of iron deficiency require immediate medical evaluation
\end{itemize}
\section{Progression and Stages}
The condition may progress through different stages of severity over time. Early stages often involve mild or subtle symptoms that may be overlooked. Moderate stages involve more noticeable symptoms affecting daily function. Advanced stages may involve significant impairment and complications.
\section{Complications}
\begin{itemize}[noitemsep]
  \item Disease progression leading to worsening symptoms
  \item Secondary infections or injuries
  \item Side effects of treatment
  \item Reduced quality of life
  \item Emotional and psychological impact
\end{itemize}
\section{Diagnosis}
Doctors diagnose this using low serum ferritin (gold standard), low serum iron, low transferrin saturation, elevated total iron-binding capacity (TIBC), elevated soluble transferrin receptor, and peripheral smear showing microcytic, hypochromic RBCs with pencil-shaped cells and target cells.

\begin{itemize}[noitemsep]
  \item Comprehensive medical history and physical examination
  \item Laboratory tests to assess organ function
  \item Imaging studies to visualize affected structures
  \item Specialized testing depending on the affected system
  \item Consultation with relevant specialists
\end{itemize}
\section{Prevention}
\begin{itemize}[noitemsep]
  \item Maintain a healthy lifestyle with balanced diet and exercise
  \item Avoid known risk factors
  \item Attend regular medical check-ups for early detection
  \item Follow recommended screening guidelines
\end{itemize}
\section{Treatment Options}
Treatment focuses on oral iron therapy (ferrous sulfate 325 mg three times daily, equivalent to 195 mg elemental iron per day) as first-line, with response monitored by reticulocytosis at 3--7 days and hemoglobin rise of 1 g/dL per 2--3 weeks

\begin{itemize}[noitemsep]
  \item Parenteral iron is used for intolerance, severe deficiency, malabsorption, or chronic kidney disease
  \item Treatment of the underlying cause is mandatory.
  \item Medications to manage symptoms and address underlying mechanisms
  \item Lifestyle modifications and self-care measures
  \item Physical therapy and rehabilitation
  \item Surgical intervention when indicated
  \item Regular monitoring and follow-up care
\end{itemize}
\section{Living with It}
\begin{itemize}[noitemsep]
  \item Follow your treatment plan as prescribed
  \item Maintain regular follow-up with your healthcare provider
  \item Make necessary lifestyle adjustments
  \item Seek support from family, friends, or support groups
  \item Stay informed about your condition
\end{itemize}
\section{Outlook}
Most people recover well with treatment.
